
\section{Poincaré 对偶}\label{sec:12-3}

为得到\BookRecordIndex{T0132}弗洛尔同调而非上同调，可以把链复形 $\operatorname{CF}_*$ 定义为函数
$\eta:\widetilde{\mathcal P}(H)\to R$ 的集合，其中函数满足有限性条件
\[
  \#\bigl\{\widetilde x\in\widetilde{\mathcal P}(H)\bigm|
  \eta(\widetilde x)\ne0,\ \mathcal A_H(\widetilde x)\ge c\bigr\}<\infty
\]
对每个 $c\in\mathbb R$ 都成立。$\Lambda_\omega$ 在 $\operatorname{CF}_*$ 上的作用由
下式给出\sourcefootnote{原书在这个公式右端把此前定义的链 $\eta$ 误印为 $\xi$；定义对象
以及紧随其后的配对\BookRecordIndex{T0348}相容式都唯一确定这里应为 $\eta(A\#\widetilde x)$。}：
\[
  (\lambda*\eta)(\widetilde x)
  :=\sum_A\lambda(-A)\eta(A\#\widetilde x).
\]
有一个显然的配对 $\operatorname{CF}^*\otimes\operatorname{CF}_*\to R$，定义为
\begin{equation}\label{eq:12-3-1}
  \langle\xi,\eta\rangle
  :=\sum_{\widetilde x}\xi(\widetilde x)\eta(\widetilde x).
\end{equation}
它与 $\Lambda_\omega$-模结构\BookRecordIndex{T0348}相容，亦即
\[
  \langle\lambda*\xi,\eta\rangle=\langle\xi,\lambda*\eta\rangle.
\]
关于这个配对，$\delta$ 的转置
$\partial:\operatorname{CF}_*\to\operatorname{CF}_*$ 为
\[
  (\partial\eta)(\widetilde y)
  =\sum_{\mu_{\mathrm{CZ}}(\widetilde x)
  =\mu_{\mathrm{CZ}}(\widetilde y)+1}
  \eta(\widetilde x)n(\widetilde x,\widetilde y).
\]
$\operatorname{CF}_*$ 的分次仍由 Conley--Zehnder 指数给出，而 $\partial$ 把次数降低一。

\BookRecordIndex{T0132}弗洛尔同调上的 PSS 同构由链映射
\[
  \Phi^\alpha:\operatorname{CM}_*(f;\Lambda_\omega)
  \longrightarrow\operatorname{CF}_*(H^\alpha),
  \qquad
  \Psi^\alpha:\operatorname{CF}_*(H^\alpha)
  \longrightarrow\operatorname{CM}_*(f;\Lambda_\omega)
\]
诱导。这些链映射按第~\ref{sec:12-1} 节中的方式定义。注意，$\Phi^\alpha$ 把次数降低
$n$，而 $\Psi^\alpha$ 把次数提高 $n$。我们选择这套记号，使得不论把 PSS 映射
$\Phi^\alpha$ 用在上同调还是同调上，它计数的底\BookRecordIndex{T0072}层几何对象都相同；有时分别把它记作
$\Phi^*_{\mathrm{PSS}}$（上同调上）和 $\Phi^{\mathrm{PSS}}_*$（同调上）。关于
$\Psi^\alpha$ 也有类似说明。

\begin{exercise}\label{ex:12-3-1}
设 $(H,J)$ 是命题~\ref{prop:12-1-2} 中的\BookRecordIndex{T0406}正则对，并定义
$\widehat H:\mathbb R/\mathbb Z\times M\to\mathbb R$ 为
\[
  \widehat H_t:=-H_{-t}.
\]
证明存在自然同构
\[
  \operatorname{HF}^k(H,J)\cong\operatorname{HF}_{-k}(\widehat H,J).
\]
若 $J$ 如注~\ref{rem:12-1-6} 那样依赖于 $t$，还须把 $J$ 换成
$\widehat J_t:=J_{-t}$。

\emph{提示：}证明映射
\[
  \widetilde{\mathcal P}(H)\longrightarrow
  \widetilde{\mathcal P}(\widehat H),
  \qquad (x,u)\longmapsto(\widehat x,\widehat u),
\]
是双射，其中 $\widehat x(t):=x(-t)$、$\widehat u(z):=u(\overline z)$；它同时反转辛作用量
和 Conley--Zehnder 指数的符号。证明 $u$ 是弗洛尔方程~\eqref{eq:12-1-3} 的解，当且仅当
$\widehat u(s,t):=u(-s,-t)$ 是把 $H$ 换成 $\widehat H$ 后方程~\eqref{eq:12-1-3}
的解。
\end{exercise}

\begin{remark}\label{rem:12-3-2}
若 $H$ 生成万有覆叠 $\widetilde{\operatorname{Ham}}(M,\omega)$ 中的元素
$\widetilde\phi$，则 $\widehat H$ 生成 $\widetilde\phi^{-1}$。更准确地说，设
$H:\mathbb R/\mathbb Z\times M\to\mathbb R$ 是含时哈密顿函数，并设
$t\mapsto\phi_t$ 是 $H$ 生成的\BookRecordIndex{T0158}哈密顿同痕。于是对每个 $t$ 都有
$\phi_{t+1}=\phi_t\circ\phi_1$。因此，在区间 $0\le t\le1$ 上，
$\operatorname{Ham}(M,\omega)$ 中的道路 $t\mapsto\phi_{-t}$ 与
$t\mapsto\phi_t^{-1}$ 保端点同伦。一个显式同伦为
\[
  \phi_t^\lambda:=\phi_{\lambda t}^{-1}\circ\phi_{(\lambda-1)t}.
\]
\end{remark}

弗洛尔上同调上的 \emph{Poincaré 对偶}同构
\[
  \operatorname{PD}^H:\operatorname{HF}^k(H)
  \overset{\cong}{\longrightarrow}\operatorname{HF}_{-k}(H)
\]
可定义为如下复合：先取习题~\ref{ex:12-3-1} 的典范等同
$\operatorname{HF}^k(H)\cong\operatorname{HF}_{-k}(\widehat H)$，再取从
$\operatorname{HF}_{-k}(\widehat H)$ 到 $\operatorname{HF}_{-k}(H)$ 的\BookRecordIndex{T0133}弗洛尔延拓映射。
在 $\operatorname{PD}^H$ 两侧复合 PSS 同构，便得到\BookRecordIndex{T0229}量子上同调上的 Poincaré 对偶同构
\[
  \operatorname{PD}:=\Psi_*^{\mathrm{PSS}}\circ
  \operatorname{PD}^H\circ\Psi^*_{\mathrm{PSS}}:
  \operatorname{QH}^k(M;\Lambda)
  \longrightarrow\operatorname{QH}_{2n-k}(M;\check\Lambda).
\]
由定理~\ref{thm:12-1-4}（再加上符号方面的考虑）可知，这个复合确实与
注~\ref{rem:11-1-20} 中的 Poincaré 对偶同构一致。于是
\[
  \langle a,\operatorname{PD}(b)\rangle
  =\langle a*b,1\rangle\in R,
  \qquad a,b\in\operatorname{QH}^*(M;\Lambda),
\]
其中左边的配对由~\eqref{eq:11-1-13} 给出（见习题~\ref{ex:11-1-22}）。

\begin{remark}[\BookRecordIndex{T0245}帽积]\label{rem:12-3-3}
固定一个上同调类 $a\in H^*(M)$。Piunikhin--Salamon--Schwarz
\srcite{323} 的一个结果断言：通过计数与 $a$ 的一个通有 Poincaré 对偶循环相交的弗洛尔
轨道而得到的弗洛尔上同调算子，等于与 $a$ 在 PSS 同构下的像作\BookRecordIndex{T0217}裤形乘积。其证明类似于
PSS 同构保持环结构的证明：计数带尖部弗洛尔圆柱，等价于计数在一个端点接上一只\BookRecordIndex{T0182}尖盘的
裤形曲面。借助 Poincaré 对偶，可把它翻译成\BookRecordIndex{T0132}弗洛尔同调算子的一个断言，而这些算子可解释为
与 $a$ 作弗洛尔\BookRecordIndex{T0245}帽积。在 $a=\operatorname{PD}([\mathrm{pt}])$ 的情形，相关断言可以改述
如下。

假定 $(H,J)$ 是\BookRecordIndex{T0132}弗洛尔同调的\BookRecordIndex{T0406}正则对。于是存在算子
\[
  \Psi_0:\operatorname{HF}_n(H,J)\longrightarrow
  \operatorname{HF}_{-n}(H,J),
\]
它通过计数方程
\[
  \partial_su+J_t(u)(\partial_tu-X_{H_t}(u))=0,
  \qquad u(s,t)=u(s,t+1),
  \qquad u(0,0)=x_0
\]
的有限能量、指标为零的解来定义；这里 $x_0\in M$ 是通有点。若
$n_H(\widetilde x,\widetilde y;x_0)$ 表示这个方程从
$\widetilde x\in\widetilde{\mathcal P}(H)$ 到
$\widetilde y\in\widetilde{\mathcal P}(H)$ 的解的带符号计数，则在链\BookRecordIndex{T0072}层次上
\[
  (\Psi_0\eta)(\widetilde y)
  :=\sum_{\mu_{\mathrm{CZ}}(\widetilde x)=n}
  \eta(\widetilde x)n_H(\widetilde x,\widetilde y;x_0).
\]
我们断言
\begin{equation}\label{eq:12-3-2}
  \Psi_0\circ\Phi_*^{\mathrm{PSS}}([M])
  =\Phi_*^{\mathrm{PSS}}([\mathrm{pt}]).
\end{equation}
由于 $1=\operatorname{PD}([M])$，这在弗洛尔上同调中等价于
\begin{equation}\label{eq:12-3-3}
  \Psi^0\circ\Phi^*_{\mathrm{PSS}}(1)
  =\Phi^*_{\mathrm{PSS}}(\operatorname{PD}([\mathrm{pt}])).
\end{equation}

有两种方法可以看出这一点。第一种方法是，从每个对 $\Psi^0$ 有贡献的弗洛尔圆柱上分离
出一只\BookRecordIndex{T0182}尖盘，由此\BookRecordIndex{T0326}推出 $\Psi^0$ 是与
$\Phi^*_{\mathrm{PSS}}(\operatorname{PD}([\mathrm{pt}]))$ 作\BookRecordIndex{T0217}裤形乘积。于是
\eqref{eq:12-3-3} 来自 $\Phi^*_{\mathrm{PSS}}$ 把\BookRecordIndex{T0228}量子杯积与\BookRecordIndex{T0217}裤形乘积交织这一事实。另一种
方法是，在每个对 $\Psi^0$ 有贡献的弗洛尔圆柱两端各粘上一只\BookRecordIndex{T0182}尖盘（如
注~\ref{rem:12-1-7}），得到带三个\BookRecordIndex{T0066}标记点的（扰动）全纯球面，并且球面在其中一个\BookRecordIndex{T0066}标记点
通过 $x_0$。于是复合
$\Psi^*_{\mathrm{PSS}}\circ\Psi^0\circ\Phi^*_{\mathrm{PSS}}$ 由与
$\operatorname{PD}([\mathrm{pt}])$ 作\BookRecordIndex{T0228}量子杯积给出。由于与 $1$ 作量子积是\BookRecordIndex{T0229}量子上同调上的
恒等映射，我们得到
\[
  \Psi^*_{\mathrm{PSS}}\circ\Psi^0\circ\Phi^*_{\mathrm{PSS}}(1)
  =\operatorname{PD}([\mathrm{pt}]),
\]
再由 $\Phi^*_{\mathrm{PSS}}=(\Psi^*_{\mathrm{PSS}})^{-1}$ 即得
\eqref{eq:12-3-3}。
\end{remark}
